            Parameter {\bf in} is the value-mode input variable or expression.
            Since it is value mode the argument may be a value mode variable,
            static mode variable or an expression containing these. An expression
            must not contain a queue read.            
            
            The module is created at the end of immediate execution from arguments
            queued by this procedure. It can be any type, primitive or compound. The
            optional output static logical variable {\bf access} goes true for one
            clock cycle to indicate that a read has occurred.
            
            This procedure can be called via the more cryptic and general
            procedure {\bf cpu\_api.read()} \autorefph{sec:lpcpuread} instead.
            
            In the case of Zynq the interface uses AXI bus GP0.            
            The size of the input data type must not exceed 256 bits as the CPU
            interface software limits burst transfers to 8 32-bit words.
            
            Boards having a serial interface to the CPU use that interface.
            The size of the input data type must not exceed 1016 bits (127 bytes).
            

            The optional output static logical variable {\bf access} goes true
            for one clock cycle to indicate that a read has occurred.
            
            The clock domain of output {\bf access} need not be the same as
            that of {\bf in}. If {\bf access} has no clock domain it will be
            given the bus clock domain.
            
            The parameters associated with this bus interface, including an
            allocated bus address, are written
            to the {\em design}.info file.

            The utility info2c is used to process this file and create
            an include file {\em design}.h which contains a funcion
            definition to access this interface. for example the 3PL
            code
            \begin{lstlisting}[gobble=12, language=threepl]
               cpu_read_value("alpha", v);
            \end{lstlisting}
            will result in info2c producing a declaration like
            \begin{lstlisting}[gobble=12, language=C]
               uint32_t alpha_read (fpga_t * fpga) {}
            \end{lstlisting}
            
            When the CPU reads a value it is sampling that value, hence data may
            be lost if the sampling rate is lower than the rate at which data
            changes. The data sampling logic ensures that the data is a valid
            sample even when the source is changing at the time of the
            read. This is a primitive interface which depends on additional
            protocols being imposed by the application code. For example to
            guarantee data integrity a write to register the data in a static
            variable prior to reading could be used. For a more complex interface see {\bf cpu\_api.read\_queue()}
            \autorefph{sec:lpcpurdqueue}.
            
            On running info2c on the .json file the following C function will have been
            defined for interface {\em name} -
            
            \begin{tabular}{| l | l |}
            \hline
            function                   &                 \\ \hline \hline
            fpga\_{\em name}\_read()   & read a datum    \\ \hline
            \end{tabular}
            
            Example: \\
            If 3PL source file prog.3pl contains the following \\
            \begin{lstlisting}[gobble=12, language=threepl]
               class(cpu, cpu_api());
               
               static(s, "uint:16");
               .
               .
               cpu.read_value("S1", s);
               .
               .
            \end{lstlisting}
            the CPU code should contain something like \\
            \begin{lstlisting}[gobble=12, language=C]
                   #include "prog_fpga.h"
                   .
                   .
                   fpga_t *fpga = NULL;
                   int     d;
                   .
                   .
                   fpga_open(portName);
                   .
                   .
                   d = fpga_S1_read(fpga);
                   .
            \end{lstlisting}
            
            This procedure creates no in-line target code so the current clock domain
            when this procedure is called is not relevant. The clock domain is not
            changed by this procedure,
            The CPU code in Python3 should contain something like \\
            \begin{lstlisting}[gobble=12, language=python]
                import sys
                sys.path.append("/Users/......./install/3pl")
                import info2py
                import prog_fpga as fpga

                fpga_open()

                fpga = info2py.info_from_file("prog")   # load info from prog.info file

                d = fpga.P1_read()

            \end{lstlisting}
